\subsubsection{文件系统入口}
在src文件夹中的main函数中实现整个文件系统的入口,实现思路为先初始化文件系统,然后在死循环中等待用户输入命令,利用之前实现的基础的文件和目录操作函数来实现复杂用户命令,具体代码如下:
\begin{ccode}
#include "fat.h"

void print_path_recursive(int dir){
    if(dir == 0){
        printf("/");
        return;
    }

    int parent = Directory[dir].parent;
    print_path_recursive(parent);

    if (parent != 0){
        printf("/");
    }

    printf("%s",Directory[dir].name);
}

void print_prompt(){
    printf("SimpleFAT:");
    print_path_recursive(current_dir);
    printf("$ ");
}

int main(){
    char line[256];
    char cmd[32];
    char arg1[64];
    char arg2[128];
    char buffer[1024];

    init_fs();

    while(1){
        print_prompt();
        if(fgets(line,sizeof(line),stdin) == NULL){
            break;
        }
        //fgets会读入\n
        line[strcspn(line,"\n")] = '\0';
        if (strlen(line) == 0){
            continue;
        }
        cmd[0] = '\0';
        arg1[0] = '\0';
        arg2[0] = '\0';

        sscanf(line, "%31s %63s %127[^\n]", cmd, arg1, arg2);
        if (strcmp(cmd, "exit") == 0) {
            break;
        }

        else if (strcmp(cmd, "ls") == 0) {
            ls();
        }

        else if (strcmp(cmd, "mkdir") == 0) {
            if (strlen(arg1) == 0) {
                printf("usage: mkdir dirname\n");
            } else if (!mkdir(arg1)) {
                printf("mkdir failed: %s\n", arg1);
            }
        }

        else if (strcmp(cmd, "cd") == 0) {
            if (strlen(arg1) == 0) {
                printf("usage: cd dirname\n");
            } else if (!cd(arg1)) {
                printf("cd failed: %s\n", arg1);
            }
        }

        else if (strcmp(cmd, "touch") == 0) {
            if (strlen(arg1) == 0) {
                printf("usage: touch filename\n");
            } else {
                fs_create(arg1);
            }
        }

        else if (strcmp(cmd, "rm") == 0) {
            if (strlen(arg1) == 0) {
                printf("usage: rm filename\n");
            } else {
                fs_delete(arg1);
            }
        }

        else if (strcmp(cmd, "open") == 0) {
            if (strlen(arg1) == 0 || strlen(arg2) == 0) {
                printf("usage: open filename mode\n");
                printf("mode: r | w | rw\n");
                continue;
            }

            int mode;

            if (strcmp(arg2, "r") == 0) {
                mode = MODE_READ;
            } else if (strcmp(arg2, "w") == 0) {
                mode = MODE_WRITE;
            } else if (strcmp(arg2, "rw") == 0) {
                mode = MODE_RDWR;
            } else {
                printf("invalid mode: %s\n", arg2);
                continue;
            }

            int fd = fs_open(arg1, mode);
            if (fd >= 0) {
                printf("opened fd = %d\n", fd);
            }
        }

        else if (strcmp(cmd, "close") == 0) {
            if (strlen(arg1) == 0) {
                printf("usage: close fd\n");
            } else {
                int fd = atoi(arg1);
                fs_close(fd);
            }
        }

        else if (strcmp(cmd, "write") == 0) {
            if (strlen(arg1) == 0 || strlen(arg2) == 0) {
                printf("usage: write fd content\n");
            } else {
                int fd = atoi(arg1);
                fs_write(fd, arg2, strlen(arg2));
            }
        }

        else if (strcmp(cmd, "read") == 0) {
            if (strlen(arg1) == 0) {
                printf("usage: read fd [size]\n");
                printf("  if size omitted, reads entire file from current offset\n");
            } else {
                int fd = atoi(arg1);

                if (fd < 0 || fd >= MAX_FD || FDTable[fd].used == 0) {
                    printf("read failed: invalid fd %d\n", fd);
                    continue;
                }

                int size;
                if (strlen(arg2) == 0) {
                    // 未指定 size，读取文件剩余全部内容
                    int dir_idx = FDTable[fd].dir_index;
                    int file_size = Directory[dir_idx].size;
                    int offset = FDTable[fd].offset;
                    size = file_size - offset;
                    if (size <= 0 || size > (int)sizeof(buffer) - 1) {
                        size = (int)sizeof(buffer) - 1;
                    }
                } else {
                    size = atoi(arg2);
                }

                int n = fs_read(fd, buffer, size);

                if (n >= 0) {
                    buffer[n] = '\0';
                    printf("%s\n", buffer);
                }
            }
        }

        else if (strcmp(cmd, "cat") == 0) {
            if (strlen(arg1) == 0) {
                printf("usage: cat filename\n");
            } else {
                int fd = fs_open(arg1, MODE_READ);
                if (fd >= 0) {
                    char cat_buf[1024];
                    int n = fs_read(fd, cat_buf, sizeof(cat_buf) - 1);
                    if (n >= 0) {
                        cat_buf[n] = '\0';
                        printf("%s\n", cat_buf);
                    }
                    fs_close(fd);
                }
            }
        }

        else if (strcmp(cmd, "seek") == 0) {
            if (strlen(arg1) == 0 || strlen(arg2) == 0) {
                printf("usage: seek fd offset\n");
            } else {
                int fd = atoi(arg1);
                int offset = atoi(arg2);

                if (!fs_seek(fd, offset)) {
                    printf("seek failed\n");
                }
            }
        }

        else if (strcmp(cmd, "cls") == 0) {
            printf("\033[2J\033[H");
            fflush(stdout);
        }

        else if (strcmp(cmd, "help") == 0) {
            printf("commands:\n");
            printf("  ls\n");
            printf("  mkdir dirname\n");
            printf("  cd dirname\n");
            printf("  touch filename\n");
            printf("  rm filename\n");
            printf("  cat filename\n");
            printf("  open filename r|w|rw\n");
            printf("  close fd\n");
            printf("  write fd content\n");
            printf("  read fd [size]\n");
            printf("  seek fd offset\n");
            printf("  cls\n");
            printf("  exit\n");
        }

        else {
            printf("unknown command: %s\n", cmd);
            printf("type 'help' for help\n");
        }
    }
    return 0;
}
\end{ccode}
其中print\_path\_recursive()函数采用递归的方式打印当前路径;print\_prompt()函数用于实现类似打印当前路径提示符;main()函数中
通过接收用户输入命令参数,在if-else分支中根据不同命令参数调用不同函数实现文件系统功能,如cat命令实现读一个文件内容并打印出来,采用的就是open+read+close的组合实现。(这里没有实现Linux中完整的cat命令所有功能,只是实现了查看指定文件名文件功能)。
